\documentclass[11pt]{article}
\usepackage[T1]{fontenc}
\usepackage[utf8]{inputenc}
\usepackage{lmodern,amsmath,amssymb,amsthm,mathtools,geometry,booktabs,microtype}
\usepackage[hidelinks]{hyperref}
\hypersetup{pdftitle={Rational Gram representations under a quadratic constraint},pdfsubject={Arithmetic classification, affine completion, and a product-cone universality boundary}}
\geometry{margin=1in}
\setlength{\emergencystretch}{2em}
\newtheorem{theorem}{Theorem}[section]
\newtheorem{lemma}[theorem]{Lemma}
\newtheorem{corollary}[theorem]{Corollary}
\newtheorem{proposition}[theorem]{Proposition}
\theoremstyle{remark}\newtheorem{remark}[theorem]{Remark}
\newcommand{\R}{\mathbb R}\newcommand{\Q}{\mathbb Q}\newcommand{\Z}{\mathbb Z}
\newcommand{\Sym}{\mathbb S}\newcommand{\PSD}{\mathbb S_+}
\newcommand{\tr}{\operatorname{tr}}\newcommand{\rank}{\operatorname{rank}}
\newcommand{\range}{\operatorname{range}}\newcommand{\spanof}{\operatorname{span}}
\newcommand{\cone}{\operatorname{cone}}\newcommand{\ri}{\operatorname{ri}}
\newcommand{\cl}{\operatorname{cl}}\newcommand{\rad}{\operatorname{rad}}
\newcommand{\diag}{\operatorname{diag}}\newcommand{\ip}[2]{\langle #1,#2\rangle}
\title{Rational Gram representations under a quadratic constraint\\\large Arithmetic local--global criteria, canonical completion domains,\\and a universality boundary for products of second-order cones}
\author{Manuscript II research continuation}
\date{15 November 2026}
\begin{document}\maketitle
\begin{abstract}
For a rational quadratic form $q(x)=x^TQx$, we characterize every rational matrix admitting a sum of rational rank-one squares whose vectors satisfy $q\geq0$. Real feasibility is a positive-semidefinite constraint and one trace inequality. Rational feasibility differs only on the trace-zero face: after restricting $q$ to the range of the target, it is equivalent to isotropy over $\Q$ of the nondegenerate quotient, unless the restriction is zero. We extend the criterion to arbitrary rational affine completion problems by a maximal-rational-support construction and at most one quadratic-form isotropy decision. This produces a canonical arithmetic column domain that is a subspace or a quadratic cone, rather than a finite polyhedral union. For the standard Lorentz cone, every real-feasible rational matrix of rank at least five factors rationally, and explicit rank-four examples show sharpness. Exact spherical second-moment problems inherit the same local--global classification. An open-domain theorem separates constrained real width from unrestricted rational Gram width. Finally, two quadratic observations on a product of $2\times2$ positive-semidefinite cones encode arbitrary rational polynomial feasibility. Thus the single-quadratic result cannot simply be iterated through products. Proofs use quadratic zero sets, rational approximation, and classical arithmetic of quadratic forms; numerical tests do not substitute for these arguments.
\end{abstract}
\paragraph{Status and attribution.} This is a complete internal proof draft, not a claim of independently certified first priority, human review, or a journal percentile. Real rank-one generation, rational orthogonal density, rational-point algorithms, the Hasse--Minkowski principle, Meyer's theorem, and small-zero bounds are explicitly credited. The new assertions isolated for comparison are the constrained arithmetic classification, its affine-carrier form, and the precise product-cone reduction. The paper does not assert a new resolution or a current undecidability status for rational polynomial feasibility.
\tableofcontents
\section{One quadratic inequality: an exact arithmetic classification}
For $Q\in\Sym^n(\Q)$ put $q(x)=x^TQx$ and
\[
 P_q=\{x\in\R^n:q(x)\geq0\},\qquad
 C_q=\cone\{xx^T:x\in P_q\}.
\]
The vector domain need not be convex. A rational permitted Gram representation is an identity $A=\sum_j u_ju_j^T$ with $u_j\in\Q^n$ and $q(u_j)\geq0$. Its width is the number of nonzero columns. Positive rational weights are equivalent to unweighted rational columns: if $a/b>0$ and $ab=\sum_{\ell=1}^4s_\ell^2$, then $(a/b)uu^T=\sum_\ell(s_\ell u/b)(s_\ell u/b)^T$. Weighted and unweighted width bounds must nevertheless be distinguished.

For $A\succeq0$, let $V=\range A$ and $r=\dim V$. If $A$ is rational, $V$ is a rational subspace. Choose any rational basis matrix $B$ of $V$. Write
\[
 A=BMB^T,\quad M\succ0,\qquad Q_V=B^TQB.
\]
Both matrices are rational. The radical of the restricted form is $R_V=\ker Q_V$ in these coordinates; its nondegenerate quotient is independent, up to rational congruence, of the chosen basis. Isotropy below means a nonzero rational zero of this \emph{nondegenerate quotient}, not merely a vector in the radical.

\begin{lemma}[Real balancing]\label{lem:balance}
For every real symmetric matrix $T$ of order $s$, there is an orthogonal $U$ with all diagonal entries of $U^TTU$ equal to $\tr(T)/s$. Consequently
\[
 C_q=\{A\succeq0:\ip{Q}{A}\geq0\},\qquad
 \cone\{xx^T:q(x)=0\}=\{A\succeq0:\ip{Q}{A}=0\}.
\]
A feasible matrix of rank $r>0$ has a real permitted factor with $r$ columns. If its trace pairing is positive, every column can have strictly positive quadratic value.
\end{lemma}
\begin{proof}
The average eigenvalue $\alpha=\tr(T)/s$ is in the Rayleigh interval. Choose a unit vector with value $\alpha$ and restrict to its orthogonal complement. The restriction has trace $(s-1)\alpha$; induction gives the claimed basis. If $A=LL^T$ is a rank-$r$ real factor, apply the assertion to $L^TQL$. Each new column of $LU$ has quadratic value $\ip{Q}{A}/r$. The reverse inclusions follow by summing quadratic values. The equality assertion uses zero average in the same proof.
\end{proof}
These real identities are instances of established rank-one generation and S-lemma phenomena; see \cite{SZ,NL}. The arithmetic assertion is different.

\begin{theorem}[Arithmetic rank-one generation]\label{thm:main}
Let $A,Q\in\Sym^n(\Q)$, $A\ne0$, and set $\sigma=\ip{Q}{A}$. A rational permitted Gram representation exists if and only if $A\succeq0$, $\sigma\geq0$, and one of the following holds:
\begin{enumerate}
\item $\sigma>0$;
\item $\sigma=0$ and $Q_V=0$;
\item $\sigma=0$, $Q_V\ne0$, and the nondegenerate quotient of $Q_V$ is isotropic over $\Q$.
\end{enumerate}
In the third case $Q_V$ is necessarily indefinite. The minimum rational permitted width in cases 1 and 2 is exactly the minimum unrestricted rational Gram width of $A$, which is at most $r+3$. In case 3 there is a representation with at most $r(r+1)/2-1$ weighted rational atoms, hence at most $4(r(r+1)/2-1)$ rational columns. The zero matrix has the empty representation.
\end{theorem}

We give the boundary proof without presuming rationality of an arbitrary real boundary factorization.
\begin{lemma}[Rational null directions]\label{lem:null}
Let $T$ be a rational indefinite symmetric matrix. If its nondegenerate quotient is anisotropic over $\Q$, its rational zeros are exactly its radical. If that quotient is isotropic, rational zeros are dense in its entire real zero cone.
\end{lemma}
\begin{proof}
Choose a rational complement to the radical. An anisotropic restriction has no nonzero rational zero, giving the first assertion. In the isotropic case choose a rational null vector $v$ not in the radical. Take rational $w$ with $v^TTw=1$, and replace it by $w-(w^TTw)v/2$. Then $q(v)=q(w)=0$ and the pairing remains one. In rational hyperbolic coordinates $x=av+bw+z$, with $z$ in their orthogonal complement, the form is
\[
 q(x)=2ab+q(z).
\]
The charts $a\ne0$ and $b\ne0$ are parametrized by rational functions:
\[
 x=av-\frac{q(z)}{2a}w+z,
 \qquad
 x=-\frac{q(z)}{2b}v+bw+z.
\]
Rational parameters are dense on each chart. A zero with $a=b=0$ is approached by $x+tv$, which is also null and is in the first chart for $t\ne0$. This proves density, including a degenerate restriction and the split binary case.
\end{proof}

\begin{lemma}[Exact rational interior reconstruction]\label{lem:reserve}
Let $\mathcal S$ be any collection of rational vectors, and suppose their rank-one atoms are dense among the atoms generating a convex cone $C$. If $C$ has rational linear span, every rational $X\in\ri C$ is a positive rational combination of at most $\dim\spanof C$ of these rational atoms.
\end{lemma}
\begin{proof}
Choose rational atoms $G_1,\ldots,G_k$ forming a basis of the span. Their sum $B$ is in $\ri C$: an open box of positive coefficients maps to a neighborhood in the span. Choose rational $\tau>0$ with $X-\tau B\in C$ and a finite real atomic representation of this residual. Approximate its atoms by rational atoms and its nonnegative weights by nonnegative rational weights, omitting any zero terms. The residual error, expressed in the rational basis $G_i$, is rational and tends to zero. Adding it to the reserve coefficients $\tau$ gives positive rational coefficients once approximation is sufficiently fine. Finally eliminate rational linear dependences while retaining nonnegative weights until at most $k$ independent atoms remain. The usual dependence elimination works whether or not the ambient cone is pointed.
\end{proof}

\begin{proof}[Proof of Theorem~\ref{thm:main}]
Every factor column lies in $V$: if $z\in\ker A$, then $0=z^TAz=\sum_j(z^Tu_j)^2$. Necessity of positive semidefiniteness and $\sigma\geq0$ follows at once.

Suppose $\sigma=0$. Every permitted column is null. Since $M\succ0$ and $\tr(Q_VM)=0$, a nonzero $Q_V$ cannot be semidefinite: conjugating by $M^{1/2}$ proves that a nonzero semidefinite restriction has a strictly signed trace. Thus it is indefinite. If its nondegenerate quotient is anisotropic, Lemma~\ref{lem:null} confines every rational column to the proper radical. Such columns cannot span $V$, contradicting $\range A=V$.

If instead the quotient is isotropic, put
\[
 C_0(V)=\{X\in\PSD(V):\ip{Q}{X}=0\}.
\]
Lemma~\ref{lem:balance} identifies this cone with the null-generated Gram cone on $V$. The point $A$ is in its relative interior: on $V$ it is positive definite, and small perturbations in the trace-zero hyperplane remain positive definite. Its span is the rational hyperplane in $\Sym(V)$ cut out by the nonzero functional $Q_V$, of dimension $r(r+1)/2-1$. Lemmas~\ref{lem:null} and \ref{lem:reserve} now give the asserted rational representation and width. If $Q_V=0$, every unrestricted rational factor of $A$ is permitted.

For $\sigma>0$, let $s$ be the minimum unrestricted rational Gram width. Such a factor exists with $s\leq r+3$ by Lemma~\ref{lem:overhead} below. Starting with any real $s$-column factor, Lemma~\ref{lem:balance} makes all column values $\sigma/s>0$. Rational factors are dense in the real Gram fiber as soon as one rational factor exists, by Theorem~\ref{thm:orbit} below. A sufficiently close rational factor preserves these strict inequalities. The resulting width is $s$, which is also the unavoidable unrestricted lower bound.
\end{proof}

\begin{corollary}[An arithmetic local--global alternative]\label{cor:localglobal}
On the nonzero trace-zero case of Theorem~\ref{thm:main}, rational permitted factorization is equivalent to isotropy of the nondegenerate quotient over every completion of $\Q$. The real condition is automatic. Failure is witnessed by a nonarchimedean place. If the nondegenerate quotient has dimension at least five, factorization always exists.
\end{corollary}
\begin{proof}
Apply the classical Hasse--Minkowski theorem, and then Meyer's theorem for indefinite rational forms in at least five variables \cite{CasselsBook}. These arithmetic theorems are used, not newly proved here.
\end{proof}

\begin{remark}[Why the nondegenerate rank is essential]
For $Q=\diag(1,-2,0,\ldots,0)$ and $A=\diag(2,1,1,\ldots,1)$, the trace is zero and $A$ has arbitrarily large rank. Nevertheless every rational null vector is in $\{x_1=x_2=0\}$, so no rational permitted factorization exists. A rational radical vector is not a successful isotropy seed. Conversely, in the split binary case $Q=\diag(1,-1)$, the target $14I_2$ needs exactly eight unweighted rational null columns: the two null lines must each contribute a sum of rational squares equal to seven, requiring four columns per line. Thus the boundary width bound is attained when $r=2$.
\end{remark}

\begin{corollary}[All fixed-range rationalizing directions]\label{cor:directions}
Let $A\ne0$ be a rational real-feasible matrix, and let $E\subseteq\Sym(V)$ be a rational linear perturbation space. The rational $D\in E$ for which $A+\varepsilon D$ factors rationally for every sufficiently small rational $\varepsilon>0$ are exactly:
\[
 \begin{cases}
 E\cap\Sym^n(\Q),&\sigma>0,\\
 \{D\in E(\Q):\ip{Q}{D}\geq0\},&\sigma=0\text{ and rational factorization holds},\\
 \{D\in E(\Q):\ip{Q}{D}>0\},&\sigma=0\text{ and rational factorization fails}.
 \end{cases}
\]
In the last case a rational rank-one direction $D=vv^T$, $v\in V$, always exists if perturbations are unrestricted within $\Sym(V)$.
\end{corollary}
\begin{proof}
For small enough $\varepsilon$, positive definiteness on $V$ and the range $V$ are preserved. Apply the classification to the trace $\sigma+\varepsilon\ip{Q}{D}$. In the last case the restricted form is indefinite, and hence has a rational strictly positive direction by continuity. Its square gives the final assertion. Thus the exceptional matrices can be rationalized without increasing their rank, but not while remaining in the trace-zero face and keeping their range.
\end{proof}

\section{Lorentz cones, a sharp rank threshold, and exact spherical moments}
Let $J_n=\diag(1,-1,\ldots,-1)$ and $\mathbb L^n=\{x:x_0\geq\sqrt{x_1^2+\cdots+x_{n-1}^2}\}$. The domain $P_{J_n}$ is the union of the two Lorentz nappes. Replacing a column by its negative does not change its square, and a nonzero null or timelike vector has nonzero time coordinate. Thus rational and real factorizations over $P_{J_n}$ and $\mathbb L^n$ are equivalent, with the same widths.

\begin{theorem}[Complete rational Lorentz criterion]\label{thm:lorentz}
Let $A\in\Sym^n(\Q)$ satisfy $A\succeq0$ and $\ip{J_n}{A}\geq0$. If the pairing is positive, $A$ has a rational Lorentz factorization. On the trace-zero face, let $V=\range A$. If the restricted form is zero, factorization also holds. Otherwise that restriction is nondegenerate of signature $(1,r-1)$, and rational Lorentz factorization is equivalent to its rational isotropy. In particular every such $A$ of rank at least five factors rationally. Failure is possible at each of ranks two, three, and four. The threshold five is optimal uniformly over matrix order.
\end{theorem}
\begin{proof}
Only the nondegeneracy assertion is additional to Theorem~\ref{thm:main}. A nonzero trace-zero restriction is indefinite, so $V$ contains a timelike vector $v$. The orthogonal complement of $v$ for the Lorentz form is negative definite. If a vector is in the radical of the restriction to $V$, it lies in this orthogonal complement and has quadratic value zero; it must therefore be zero. There can be at most one positive direction. Apply Corollary~\ref{cor:localglobal}.
\end{proof}

\paragraph{Rank two.}
Take
\[
 A_2=\begin{pmatrix}2&0&0\\0&1&1\\0&1&1\end{pmatrix}.
\]
Its range is $\{(t,u,u)\}$, and the restricted Lorentz form is $t^2-2u^2$. It has no nonzero rational zero. Nevertheless $A_2\succeq0$ and $\ip{J_3}{A_2}=0$, so it is real-feasible. Its two real columns can be $(1,\pm1/\sqrt2,\pm1/\sqrt2)^T$.

\paragraph{Rank three.}
Let $B_3(t,u,v)=(t,u,u,v,v,v)$ and set $A_3=B_3\diag(5,1,1)B_3^T$. The restricted form is $t^2-2u^2-3v^2$. If an integral primitive vector were null, reduction modulo three would force $3\mid t,u$. The equation would then force $3\mid v$, a contradiction. Thus no rational permitted factor exists. An explicit real representation uses the four coefficient columns $(1,\pm1/\sqrt2,0)$ and $(\sqrt{3/2},0,\pm1/\sqrt2)$, mapped by $B_3$.

\paragraph{Rank four.}
The particularly small example
\[
 A_4=\begin{pmatrix}
48&36&0&0&0\\36&27&0&0&0\\
0&0&7&0&0\\0&0&0&7&0\\0&0&0&0&7
\end{pmatrix}
\tag{R4}\label{eq:r4}
\]
has rank four and $\ip{J_5}{A_4}=0$. A rational range basis maps $(t,u,v,w)$ to $(4t,3t,u,v,w)$, so the restricted form is $7t^2-u^2-v^2-w^2$. A primitive integral zero cannot have odd $t$, since three squares cannot sum to seven modulo eight. If $t$ is even, reduction modulo four forces $u,v,w$ even as well, again contradicting primitiveness. It has no rational nonzero null vector. The six real columns $(2\sqrt2,3/\sqrt2,\pm\sqrt{7/2}e_i)$, $i=1,2,3$, do give the matrix. None of these examples contradicts the orthant theorems: their individual columns are constrained to a different cone.

\begin{corollary}[Totally-real descent already fails on a curved cone]
The rank-two matrix $A_2$ above has a Lorentz factorization over the totally real field $\Q(\sqrt2)$ but has no rational Lorentz factorization of any finite width.
\end{corollary}
\begin{proof}
Its two displayed real columns belong to $\Q(\sqrt2)^3$. The field has two real embeddings, so it is totally real. The anisotropic range restriction already excludes all rational factorizations.
\end{proof}
Thus the earlier polyhedral constrained totally-real descent theorem cannot be extended to curved domains without a new hypothesis. The arithmetic local--global condition on the restricted quadratic form is the corrected criterion; the name or signature of a factor field alone is not sufficient.

\subsection{An exact second-moment classification}
For rational $\mu\in\Q^k$ and $S\in\Sym^k(\Q)$ put
\[
 M=\begin{pmatrix}1&\mu^T\\\mu&S\end{pmatrix}.
\]
A finite probability law supported on the unit ball, with first moment $\mu$ and second moment $S$, exists exactly when $M\succeq0$ and $\tr S\leq1$. For the unit sphere replace the last inequality by equality. This follows from Lemma~\ref{lem:balance}: orient each nonzero Lorentz column $(t,z)$ with $t>0$, and assign probability $t^2$ to $z/t$. Conversely a finite law is a Gram representation using columns $\sqrt p(1,z)$. The same equivalence holds for rational support and rational probabilities versus weighted rational Gram representations. Unweighted rational columns follow from four squares if needed.

\begin{corollary}[Rational spherical moments]\label{cor:sphere}
Assume $M\succeq0$ and $\tr S=1$. The prescribed moments have a finite realization on rational sphere points with rational probabilities if and only if the restricted Lorentz form on $\range M$ is zero or is isotropic over $\Q$. They always have such a realization when $\rank M\geq5$. For $\tr S<1$ there is always a rational realization supported strictly inside the ball.
\end{corollary}
The rank-four obstruction is geometrically transparent. Divide \eqref{eq:r4} by $48$. It prescribes first moment $(3/4,0,0,0)$ and second moments $\diag(9/16,7/48,7/48,7/48)$. Zero variance in the first coordinate forces every atom to have that coordinate $3/4$. The remaining coordinates must satisfy
\[
 z_2^2+z_3^2+z_4^2=7/16.
\]
There is a real six-atom law, but no rational point on this required latitude. Thus rational moment data and real feasibility do not by themselves give rational atoms, even on a smooth sphere.

The same theorem covers an arbitrary rational ellipsoid described by a single quadratic inequality. Homogenize its defining polynomial using the first coordinate and apply Theorem~\ref{thm:main}; no rational square root or rational congruence to a standard ball is needed. More generally it applies to nonconvex regions described by one quadratic inequality, with homogeneous lifts when appropriate. An affine polynomial support problem must retain the positive homogenizing coordinate: in the ball and ellipsoid cases it is automatic for nonzero permitted columns; it is not automatic for arbitrary unbounded quadrics.

\section{A canonical arithmetic carrier for every rational affine completion}
We now allow the target matrix itself to be unknown. Let $H\subseteq\Sym^n$ be rational affine, and write
\[
 \mathcal F=H\cap\PSD^n\cap\{X:\ip{Q}{X}\geq0\}.
\]
The following construction needs a decision procedure for rational points in convex semialgebraic sets. Such a procedure, including strict inequalities, is a classical algorithmic primitive \cite{SEDZ}; it must not be confused with a general rational-point procedure for nonconvex varieties.

\begin{lemma}[Maximal rational support]\label{lem:maxsupport}
There is an effective procedure that either finds $\mathcal F\cap\Sym^n(\Q)=\varnothing$, or returns a rational $X_*\in\mathcal F$ with
\[
 S=\range X_*=\sum_{X\in\mathcal F\cap\Sym^n(\Q)}\range X.
\]
It uses at most $n+1$ rational-point queries. The same assertion holds when the inequality is omitted or the allowed matrix support is restricted to a rational subspace.
\end{lemma}
\begin{proof}
Find a rational feasible matrix. If its range is $S$, form the rational Euclidean projector $\Pi$ onto $S^\perp$ and ask whether the convex semialgebraic set $\mathcal F\cap\{X:\tr(\Pi X)>0\}$ has a rational point. For a positive-semidefinite matrix, $\tr(\Pi X)=0$ is equivalent to $\range X\subseteq S$: factor $X$ and sum the squared projected column norms. When $S$ is the whole allowed ambient subspace, stop without a redundant query. A new point increases the range after averaging it with the previous point. Both feasibility and rationality are preserved, and the range of a sum of positive-semidefinite matrices is the sum of their ranges. There are at most $n$ increases, including the initial positive-rank step if one occurs.
\end{proof}

For a rational permitted column $u$, call it viable in $H$ when some rational permitted Gram matrix $X\in H$ has $X-\tau uu^T$ also rationally permitted for some rational $\tau>0$.

\begin{theorem}[Quadratic arithmetic carrier]\label{thm:carrier}
For every rational $Q,H$, the following algorithm is finite and computes either an empty rational Gram fiber or its unique closed arithmetic column carrier $P_*$. In the nonempty case
\begin{align*}
 \{\text{viable rational columns}\}&=P_*\cap\Q^n,\\
 \{\text{rational permitted Gram matrices in }H\}
 &=\{\text{rational }P_*\text{-Gram matrices in }H\},\\
 \cl\{\text{rational permitted Gram matrices in }H\}&=C(P_*)\cap H.
\end{align*}
Moreover $H\cap\ri C(P_*)\ne\varnothing$, and one rational matrix there witnesses every viable rational column. The carrier is one of a rational subspace, a single quadratic cone inside a rational subspace, or its quadratic zero cone.

The algorithm makes at most $2n+3$ rational convex-semialgebraic point queries and at most one rational quadratic-form isotropy decision.
\end{theorem}
\begin{proof}
First compute the maximal rational support $S$ of $\mathcal F$ and a witness $X_*$. If there is no rational feasible matrix, there can be no rational factorization. If $S=\{0\}$, return carrier $\{0\}$, with the empty factorization.

Ask whether there is a rational matrix in $\mathcal F$ with positive trace pairing. If yes, average it with $X_*$ so that $X_*$ has both range $S$ and positive pairing. Set
\[
 P_*=S\cap P_q.\tag{C1}
\]
Every rational feasible Gram column lies in $S$. Conversely, for any rational $u\in P_*$, a sufficiently small rational subtraction $X_*-\tau uu^T$ remains positive definite on $S$ and has positive trace pairing. Theorem~\ref{thm:main} supplies a rational representation of the residual. Thus every such column is viable.

Otherwise every rational matrix in $\mathcal F$ has trace pairing zero. The restriction $q_S$ is either identically zero or indefinite, by the positive definiteness of $X_*$ on $S$. If it is zero, set
\[
 P_*=S.\tag{C2}
\]
If it is nonzero and its nondegenerate quotient is rationally isotropic, set
\[
 P_*=\{u\in S:q(u)=0\}.\tag{C3}
\]
In either case Theorem~\ref{thm:main} shows that $X_*$ and all sufficiently small rational residuals $X_*-\tau uu^T$ have permitted factors. Conversely every column in any rational permitted feasible representation is a rational null vector in $S$.

In the remaining case, the rational null vectors of $q_S$ are exactly its radical $T$. Every rational permitted completion must therefore be supported in $T$. Apply Lemma~\ref{lem:maxsupport} to $H\cap\PSD(T)$. If no rational matrix exists, the rational Gram fiber is empty. Otherwise let $S_0\subseteq T$ be its maximal rational support, replace $X_*$ by that witness, and set
\[
 P_*=S_0.\tag{C4}
\]
The form is zero on $S_0$, so unrestricted rational PSD factors suffice. This completes the algorithm; there is no second isotropy test.

In (C1), rational points are dense in $P_*$: the restricted form has a positive direction, null points are limits of strictly positive points, and rationals are dense in that open set. In (C3) density is Lemma~\ref{lem:null}; in (C2) and (C4) it is immediate. We have already proved exact equality of rational viable directions, so taking their closure proves that $P_*$ is canonical.

In every case $X_*$ is in $\ri C(P_*)$: it is positive definite on the carrier span, with either strictly positive trace or the appropriate trace equality. The real cone $C(P_*)$ is closed, either by Lemma~\ref{lem:balance} or as a PSD face. Thus limits of rational feasible factors lie in its slice. Conversely, take $Y\in C(P_*)\cap H$. The segment $(1-t)Y+tX_*$, $t>0$, lies in the relative interior. In the rational affine space $H\cap\spanof C(P_*)$, approximate its points by rational points still in that relative interior. Theorem~\ref{thm:main} gives permitted rational factors for these points in cases (C1)--(C4). Letting the approximation error and then $t$ tend to zero proves the closure formula. This proof makes no claim that a nonclosed linear image is closed.
\end{proof}

\begin{corollary}[Exact observation-preserving regularization]
Let $A\in H\cap C_q$ be rational, and compute $P_*$ from Theorem~\ref{thm:carrier}. There is rational $D$ with $\mathcal L(D)=0$ and rational permitted factorizations of $A+tD$ for every rational $0<t\leq1$ if and only if
\[
 A\in C(P_*).
\]
For sufficiency take $D=X_*-A$. Necessity follows from the closure identity. The result permits rank changes and is therefore stronger than Corollary~\ref{cor:directions}.
\end{corollary}

For the positively oriented Lorentz cone, intersect the computed carrier with the positive nappe. Sign changes preserve rank-one squares, so the proof and completion criterion remain valid.

The carrier is not necessarily polyhedral. For $Q=\diag(1,1,-1)$ and $A=\diag(1,1,2)$, the full-matrix fiber has carrier $\{u:u_1^2+u_2^2-u_3^2=0\}$. A convex cone contained in one nappe of this null cone has only one ray; the null cone has infinitely many rays. Therefore no finite union of convex polyhedral cones can be its carrier. The earlier finite-polyhedral carrier conclusion does not extend literally to curved domains; the quadratic carrier above is its corrected replacement.

\paragraph{Exact failure certificates.}
If a real feasible fiber contains no rational matrix, the rational-point procedure supplies the relevant exact arithmetic decision, not a false real separator. If it contains rational matrices but the isotropy test fails, the certificate additionally identifies the maximal rational support and the radical restriction. The Hasse--Minkowski obstruction is a local arithmetic certificate, not a separating hyperplane from a real-feasible cone. Weak infeasibility of a measurement image remains possible and does not justify a strict real separator.

\section{Open constraints: an exact width law and quantitative rational density}
This section explains why the arithmetic obstruction in the main theorem is concentrated on quadratic zero sets. No polyhedral assumption is needed for open column constraints.

\begin{theorem}[Rational Gram-fiber density]\label{thm:orbit}
For rational $A\succeq0$ and an integer $s$, let
\[
 \mathcal M_s(A)=\{B\in\R^{n\times s}:BB^T=A\}.
\]
Either this set has no rational point, or its rational points are dense in the whole real set, including rank-deficient fibers.
\end{theorem}
\begin{proof}
Fix a rational point $B_0$. If $B\in\mathcal M_s(A)$, the correspondence between the rows of $B_0$ and the rows of $B$ is an isometry on their row spans, because their row Gram matrices agree. It is well defined even when rows are dependent: a linear relation has squared norm zero on both sides. Extend this isometry to $\R^s$ to obtain $B=B_0O$ with $O$ orthogonal. Rational orthogonal matrices are dense, as the explicit signed Cayley construction below proves. Thus $B_0O_j\to B$ with every $B_0O_j$ rational and its Gram matrix exactly $A$.
\end{proof}
Rational density of the orthogonal group and signed Cayley charts are classical \cite{Schmutz,TWZ}; the exact constrained-width consequence is separated from those antecedents.

\begin{theorem}[Exact open-domain width]\label{thm:width}
Let $A\ne0$ be rational PSD, $V=\range A$, and let $P\subseteq V$ be relatively open and closed under multiplication by positive scalars. Suppose $A$ has a real Gram representation with columns in $P$. Write $s_\R(P,A)$ for its least real permitted width and $s_\Q(A)$ for its least unrestricted rational Gram width. Then the least rational permitted width is
\[
 \boxed{s_\Q(P,A)=\max\{s_\R(P,A),s_\Q(A)\}.}
\]
Here $s_\Q(A)\leq\rank A+3$ and $s_\R(P,A)\leq r(r+1)/2$ for $r=\rank A$.
\end{theorem}
\begin{proof}
The two lower bounds are immediate. Let $s$ be their maximum. Split any nonzero real column into sufficiently many equal positively rescaled copies to obtain a real permitted factor of exactly $s$ columns. Pad an unrestricted rational factor with zero columns to width $s$. Theorem~\ref{thm:orbit} gives rational factors of that width arbitrarily close to the real permitted factor; openness preserves all column constraints. Lemma~\ref{lem:overhead} gives $s_\Q(A)\leq r+3$. Conic Carath\'eodory gives the real-width bound, with positive weights absorbed into columns.
\end{proof}
\begin{lemma}[Sharp classical rational-width overhead]\label{lem:overhead}
Every rational PSD matrix of rank $r>0$ has an unrestricted rational Gram representation of width at most $r+3$. This bound is sharp for every $r$.
\end{lemma}
\begin{proof}
Rational diagonalization reduces the assertion to representing the positive diagonal form $\diag(a_1,\ldots,a_r)$. Start in rational Euclidean space of dimension $r+3$. After choosing $i-1$ mutually orthogonal rational rows with prescribed squared lengths, their rational orthogonal complement has dimension $r+4-i\geq4$ and a rational positive-definite form $g$. The form $g(x)-a_it^2$ is indefinite in at least five variables, so Meyer's theorem supplies a nonzero rational zero. Positive definiteness makes $t\ne0$. The vector $x/t$ in that complement has squared length $a_i$. Induction constructs the rows.

For sharpness use $\diag(I_{r-1},7)$. Any rational factor with at most $r+2$ columns would give $r-1$ rational orthonormal rows followed by a rational orthogonal row of squared length seven. Rational Householder reflections carry the orthonormal rows to standard coordinate rows: reflecting a rational unit vector $v\ne e_1$ in the rational vector $v-e_1$ sends it to $e_1$, and induction acts on the remaining coordinates. The last row would express seven as at most three rational squares. After clearing denominators, $u^2+v^2+w^2=7t^2$ has no primitive integral solution: odd $t$ is impossible modulo eight, while even $t$ forces $u,v,w$ even modulo four. On the other hand, $7=2^2+1^2+1^2+1^2$ supplies $r+3$ columns. This is a consequence of classical quadratic-form arithmetic, not a claim of a new unrestricted width theorem.
\end{proof}
For $A=[7]$ and $P=(0,\infty)$, the real permitted width is one and the rational permitted width is four. Thus even an open scalar constraint cannot eliminate the arithmetic width term.

\subsection{Explicit precision and height of the orthogonal approximation}
For an orthogonal $O$ of order $s$, average $\det(I+DO)$ over the $2^s$ diagonal sign matrices $D$. Multilinearity in the rows cancels every nonconstant sign monomial, so the average is one. Choose $D$ with determinant at least one. The matrix $I+DO$ has norm at most two, and hence its smallest singular value is at least $2^{-(s-1)}$. Put
\[
 U=DO,\qquad T=(I-U)(I+U)^{-1}.
\]
Then $T^T=-T$ and $\|T\|_2\leq2^s$. Round its strict upper triangle to multiples of $2^{-h}$ and use skew symmetry to define $T_h$. Then
\[
 \|T-T_h\|_F\leq s2^{-h}.
\]
For every real skew-symmetric $T$, $\|(I+T)^{-1}\|_2\leq1$. The resolvent identity for
\[
 \mathcal C(T)=(I-T)(I+T)^{-1}=2(I+T)^{-1}-I
\]
gives
\[
 \|D\mathcal C(T_h)-O\|_F\leq2s2^{-h}.
\]
The approximant is exactly rational and orthogonal, not approximately orthogonal.

Call $H$ a coordinate height bound if every reduced numerator and positive denominator has magnitude at most $2^H$. The integer matrix $2^h(I+T_h)$ has entries of magnitude at most $2^{h+s+1}$. Determinant and cofactor bounds therefore give the conservative common-denominator height
\[
 H_O=3+s(h+s+1)+\lceil\log_2(s!)\rceil.
\]
If the rational base factor $B_0$ has coordinate height $L_0$, every entry of $B_0D\mathcal C(T_h)$ has height at most
\[
 sL_0+H_O+\lceil\log_2 s\rceil+2.
\]
If a real permitted factor $B_0O$ has distance at least $\delta>0$ from the complement of the permitted subset in its Gram fiber's ambient column space, choose
\[
 h\geq\max\left\{0,\left\lceil\log_2\frac{4s\max(1,\|B_0\|_2)}{\delta}\right\rceil\right\}.
\]
The rational factor is then within $\delta/2$. This is a height theorem conditional on a base rational factor and a clearance bound; it is not a polynomial discovery theorem for arbitrary open sets.

\subsection{The same zero-set principle gives a quantitative dual margin}
Suppose $A=\sum_j u_ju_j^T$ has nonzero columns in $P\subseteq V$, and for some $0<\theta\leq1/2$ the relative ball
\[
 \{v\in V:\|v-u_j\|_2<\theta\|u_j\|_2\}
\]
is contained in $P$ for every column. Set $\lambda=\lambda_{\min}(A|_V)>0$. Then
\begin{equation}\label{eq:openball}
 \{A+E:E\in\Sym(V),\ \|E\|_F\leq\theta\lambda/2\}\subseteq C(P).
\end{equation}
Indeed, with square roots and inverses on $V$, take
\[
 T=(A+E)^{1/2}A^{-1/2},\qquad \widetilde u_j=Tu_j.
\]
Their Gram sum is $A+E$. The square-root difference solves the Sylvester equation
\[
 (A+E)^{1/2}Z+ZA^{1/2}=E,
 \qquad Z=(A+E)^{1/2}-A^{1/2}.
\]
Its integral inverse, or diagonalization on the two sides, gives
\[
 \|Z\|_F\leq\frac{\|E\|_F}{\sqrt{\lambda-\|E\|_2}+\sqrt\lambda}.
\]
Consequently $\|T-I\|_2\leq\|E\|_F/\lambda\leq\theta/2$, proving that each new column is still permitted. The radius in \eqref{eq:openball} is not asserted optimal. It implies the dual estimate
\[
 \ip{W}{A}\geq\frac{\theta\lambda}{2}\|W|_V\|_F
 \quad\text{whenever }u^TWu\geq0\text{ on }P.
\]
Here $W|_V$ denotes Euclidean orthogonal compression. Thus openness gives both exact rational persistence and a quantitative positive dual margin.

Qualitatively the mechanism is especially simple. If a nonnegative quadratic form vanishes at an interior permitted column, two-sided variations force its restriction to $V$ to be PSD and to annihilate the column. If it vanishes on all columns of a factor spanning $V$, the restriction is zero. Arithmetic difficulty begins when a supporting quadratic zero set is curved or otherwise lacks these two-sided variations. The single-quadratic criterion identifies that difficulty exactly by isotropy.

\section{A universality boundary: products of second-order cones}
Generalized copositive representations of nonconvex quadratic programs over products of second-order cones and over PSD cones are already established in Burer--Dong \cite{BD}. Their real convexification result is a direct antecedent, not a new claim here. Our statement below isolates two homogeneous observations, exactly preserves rational solutions without a boundedness promise, and excludes any escape through factor width.

The preceding theorem handles every one-quadratic column inequality. It cannot be extended to arbitrary products by treating the factors separately: simultaneous boundary conditions can encode arbitrary rational polynomial equations. The following is an unconditional reduction, independent of any current decision-theoretic status assigned to that equation problem.

\begin{theorem}[Two-observation universality]\label{thm:universal}
Given finitely many rational polynomial equations, represented by arithmetic circuits, one can construct in polynomial time a product cone
\[
 P=(\PSD^2)^k
\]
and two rational homogeneous quadratic forms $F,m$ on its vector space such that the following are equivalent:
\begin{enumerate}
\item the polynomial system has a rational solution;
\item there are finitely many rational $Z^{(a)}\in P$ with
\[
 \sum_a F(Z^{(a)})=0,\qquad \sum_a m(Z^{(a)})=1;
\]
\item one rational $Z\in P$ satisfies $F(Z)=0$, $m(Z)=1$.
\end{enumerate}
The same equivalence holds over $\R$. In particular, allowing an arbitrary factor width does not remove the encoded arithmetic condition.
\end{theorem}
\begin{proof}
Convert the circuits to rational affine-linear gates and squaring gates. Multiplication uses
\[
 ab=((a+b)^2-(a-b)^2)/4,
\]
so the conversion increases size only by a constant factor. Introduce one PSD block for every scalar register:
\[
 Z_j=\begin{pmatrix}t_j&u_j\\u_j&v_j\end{pmatrix}.
\]
Use homogeneous linear equations to equate all $t_j$ to a common $t=t_1$, to encode every affine-linear gate after homogenization, and to identify $u_{\rm out}=v_{\rm in}$ for every squaring gate. A constant register $c$ has equation $u_j=ct$. Output equations are $u_{\rm out}=0$. Denote all these homogeneous linear residuals by $\ell_\nu$ and set
\[
 F(Z)=\sum_{j=1}^k(t_jv_j-u_j^2)+\sum_\nu\ell_\nu(Z)^2,
 \qquad m(Z)=t_1^2.
\]
Every summand in $F$ is nonnegative on the product PSD cone. A feasible sum of atoms therefore forces every determinant and every linear residual to vanish in each atom. Some atom has $t_1\ne0$ because the sum of its squares is one. Positive semidefiniteness and the common-time equations give $t>0$. Divide the chosen atom by $t$. Its registers $a_j=u_j/t$ obey
\[
 v_j/t=(u_j/t)^2=a_j^2,
\]
and all circuit equations. They supply a rational solution.

Conversely, from a rational circuit solution put $t_j=1$, $u_j=a_j$, and $v_j=a_j^2$. All blocks are rational PSD of rank one, and this single atom has $(F,m)=(0,1)$. The argument over the reals is identical. All matrices and coefficients have descriptions of polynomial size in the circuit description.
\end{proof}
The product cone is a product of rationally equivalent Lorentz cones: the linear map
\[
 \begin{pmatrix}a&b\\b&c\end{pmatrix}
 \longmapsto ((a+c)/2,(a-c)/2,b)
\]
maps $\PSD^2$ onto $\mathbb L^3$. The observation construction is a genuine matrix completion problem: after vectorizing the symmetric blocks, $F(Z)$ and $m(Z)$ are linear functions of the rank-one matrix $\operatorname{vec}(Z)\operatorname{vec}(Z)^T$. Thus the theorem concerns exactly two specified linear observations of a column-constrained Gram matrix.

\begin{corollary}[A single PSD column domain also contains the reduction]
Theorem~\ref{thm:universal} remains true with the domain replaced by one cone $\PSD^{2k}$ of permitted matrix-valued columns, and the same two-observation target.
\end{corollary}
\begin{proof}
Partition a symmetric matrix into $k$ principal blocks of order two. Add to $F$ the sum of squares of all off-block entries, and use the preceding formula on the diagonal blocks. Each principal block is PSD, so all determinant terms are nonnegative. A zero of $F$ is block diagonal and is exactly a product-cone zero. Vectorizing the single PSD matrix again turns both quadratic forms into linear observations of its rank-one Gram square.
\end{proof}

\paragraph{A rational full matrix with only irrational permitted columns.}
The circuit $x^2=2$ gives a compact exact example in the four-block domain. In the order $(t_j,u_j,v_j)$ for the registers $x,x^2,2,x^2-2$, put
\[
 a=(1,0,2,\ 1,2,4,\ 1,2,4,\ 1,0,0)^T,\qquad b=e_2,
 \qquad A=aa^T+2bb^T.
\]
The vectors $a\pm\sqrt2b$ encode the two real circuit solutions and belong to the product PSD domain. Therefore
\[
 A=\frac12(a+\sqrt2b)(a+\sqrt2b)^T+
   \frac12(a-\sqrt2b)(a-\sqrt2b)^T
\]
is a rational matrix in the real constrained Gram cone. Its two observations are $(0,1)$. Any rational permitted factorization of this \emph{full matrix}, of any width, would supply a rational solution of $x^2=2$ by Theorem~\ref{thm:universal}. None exists. The real permitted factor can be taken over the totally real field $\Q(\sqrt2)$ by absorbing $1/2$ into its columns. This is a separate illustration of the failure of unqualified constrained totally-real descent outside the polyhedral setting.

\begin{corollary}[Unconditional output-height obstruction in variable dimension]
For a circuit with $s$ successive squarings, there are two-observation product-cone instances of description length $O(s\log(s+2))$ such that every rational factor certificate contains an entry of coordinate height at least $2^{s-1}$, regardless of its width.
\end{corollary}
\begin{proof}
Take the system $a_0=2$, $a_{j+1}=a_j^2$ for $j=0,\ldots,s-1$. Every positive-time atom satisfies $u_s/t=2^{2^s}$. If all its rational entries had numerator and denominator bounded by $2^H$, this ratio would be at most $2^{2H}$. Hence $H\geq2^{s-1}$. At least one positive-time atom is unavoidable. The circuit and block descriptions use only small coefficients; indices account for the stated bit length. This obstruction grows the dimension and does not contradict the fixed-dimensional bounds for one quadratic constraint.
\end{proof}

\paragraph{What the reduction does and does not say.}
It proves that a total decision algorithm for all these product-cone completion problems would decide arbitrary rational polynomial feasibility, and a uniform computable height bound would permit bounded search for that problem. It does not assert that such an algorithm is impossible, nor does it assume an unresolved or undecidable status for rational polynomial feasibility. Any claim about that status would need an independent current literature audit. The equivalence, rather than a historical status assertion, is the mathematical obstruction supplied here.

The distinction is structural. Polyhedral zero domains linearize into rational coefficient kernels. A single quadratic zero domain has rational parametrizations precisely when its nondegenerate part is isotropic. Intersections of independent curved boundary conditions can implement arithmetic circuits, even though each individual block is only a three-dimensional second-order cone. This is why a universal factorization theorem must specify the geometry of its zero domains rather than only the closedness, rational description, or interior conditioning of the original cone.

\section{Effective constructions, bit resources, and scope}
We separate mathematical existence, output height, decision cost, and implementations. Let $L$ denote total binary input length, including $Q$, the rational matrix or affine equations, and the ambient dimension. All constants indexed by $n$ below are effective. No polynomial dependence on variable dimension is asserted.

\subsection{Rational support computation}
Positive semidefiniteness can be expressed by $2^n-1$ principal-minor inequalities, of degree at most $n$, in $d=n(n+1)/2$ matrix coordinates. Redundant affine equations can be removed. Safey El Din--Zhi's rational-point algorithm for convex semialgebraic sets \cite[Theorem 1.1]{SEDZ} has bit cost
\[
 \sigma^{O(1)}(sD)^{O(k^3)}
\]
for $s$ predicates of degree at most $D$ in $k$ variables and coefficient height $\sigma$, with output height $\sigma D^{O(k^3)}$. Strict inequalities are included in its quantifier-free input model. In our queries $k\leq d$, $D\leq n$, and $s\leq2^n+O(d)$. Rational projections and averaging use only matrices of fixed order and increase heights by a dimension-dependent factor. Thus at most $2n+3$ queries, together with the linear algebra between them, cost $\operatorname{poly}_n(L)$ bit operations and have output coordinate height $O_n(L)$. This is a known primitive; it is not a new general rational-point algorithm.

\subsection{The isotropy resource is not a free oracle}
For a rational nondegenerate quotient of order $k$, clear denominators to obtain an integral form whose coefficient magnitudes are at most $M=2^{O_n(L)}$. Cassels' classical small-zero theorem, in the explicit form recorded in \cite[Theorem 1.1]{Fukshansky}, supplies
\[
 0<\|v\|_\infty\leq
 \left(3\sum_{i,j}|f_{ij}|\right)^{(k-1)/2}
 \leq(3k^2M)^{(k-1)/2}
\]
whenever the integral form $v^TFv$, $F=(f_{ij})$, has a rational zero. See also \cite{CasselsSmall,CasselsBook,Dietmann}. Enumerating the finite box up to the ceiling of this bound is a total decision-and-construction procedure: every test is an exact integer evaluation. Its conservative fixed-dimensional bit cost is $2^{O_n(L+1)}$, and a successful seed has coordinate height $O_n(L+1)$. Classical local-symbol algorithms can be used instead, but their arithmetic subroutines, including any factoring assumptions, must be charged. We do not claim an unqualified deterministic polynomial-time isotropy algorithm.

The classification therefore has at most one isotropy decision in addition to the rational support queries. Once a positive answer is accompanied by a rational isotropic seed of height $h$, the remaining weighted-factor construction described next is polynomial in $L+h$ for fixed $n$. Since such a seed of height $O_n(L)$ exists, every positive instance has a polynomial-height certificate, independently of the chosen isotropy algorithm's running time.

\subsection{Constructing the exact factors}
For clarity we expand the reserve argument into effective finite steps.

In the isotropic boundary case, choose the rational hyperbolic basis of Lemma~\ref{lem:null}. In each of its charts with $a=1$ or $b=1$, the vector coordinates are polynomials of degree at most two in $z$, and the outer-product coordinates have degree at most four. Evaluate on the tensor set $\{0,1,2,3,4\}^{r-2}$. Multivariate interpolation shows that these finitely many rational null atoms span the same matrix space as the complete charts. The remaining null atoms are their limits up to scaling, so their squares lie in that span as well. Select a rational atom basis $G_i$ by elimination and put $B=\sum_iG_i$. All basis and coefficient heights are $O_n(L+h)$. Since their span contains the positive-definite-on-$V$ target, $B$ is also positive definite on $V$.

Select a positive rational $\tau$ of height $O_n(L+h)$ such that $A-\tau B\succeq0$ on $V$. This is possible from the determinant and norm bounds for rational positive-definite matrices of fixed order; alternatively successively test dyadic $\tau$ using rational principal minors. The trace condition is preserved. Lemma~\ref{lem:balance} gives a real null factor of the residual. Standard real-algebraic sampling of the equations $CC^T=A-\tau B$, $q(C_j)=0$ uses a number of variables and degree bounded in terms of $n$ only, and provides algebraic columns of bounded degree and height $O_n(L+h)$ \cite{SEDZ}. Approximate them within the rational hyperbolic charts. A nonzero algebraic chart coordinate has a root-separation lower bound $2^{-O_n(L+h)}$; if both hyperbolic coordinates vanish, first perturb by a small rational multiple of $v$ as in Lemma~\ref{lem:null}. Thus the rational parameters and atom approximations needed for accuracy $2^{-p}$ have height $O_n(L+h+p)$.

The inverse matrix of the selected rational atom basis has height $O_n(L+h)$ by Cramer's rule. Taking $p=O_n(L+h)$ therefore makes the residual basis-coefficient error smaller than $\tau/2$. All corrected coefficients are positive, the trace equality remains exact, and the reconstructed matrix is exactly $A$. Rational dependence elimination reduces to at most $r(r+1)/2-1$ weighted atoms.

In the positive-trace case, find a rational vector with $q>0$ by rational diagonalization. A rational box around it remains in $q>0$, with radius of height $O_n(L)$ obtainable by direct quadratic estimates. Evaluating rank-one squares on a degree-two interpolation grid in this box gives a rational atom basis of $\Sym(V)$. Choose $\tau$ so that $A-\tau B$ stays positive definite on $V$ and has positive trace. Sample a real factor with strictly positive column values; fixed-dimensional real-algebraic height bounds give sufficient positive clearance. Rational rounding and the same reserve repair produce at most $r(r+1)/2$ weighted atoms of height $O_n(L)$. If $q|_V=0$, ordinary rational diagonalization suffices. These constructions also apply to the rational reserve matrix returned by the affine-carrier algorithm.

A positive rational coefficient $a/b$ can be converted to at most four rational squares with coordinate height $O(\operatorname{bits}(a)+\operatorname{bits}(b))$. This is an existence bound. For a deterministic polynomial-time implementation avoiding a four-square search, expand the integer $ab$ in binary: an even power of two is a square, and an odd power is a sum of two equal squares. Dividing by $b$ gives $O(\operatorname{bits}(ab))$ rational squares. Therefore deterministic conversion may output $O_n(L+h)$ columns rather than the four-times-weighted-width existence bound.

\subsection{Finite precision and observations}
A rational target in the terminal slice can be approximated within $2^{-p}$ with its affine measurements unchanged by moving a rational distance toward $X_*$ and applying the preceding construction. Fixed-dimensional determinant estimates control the new relative-interior margin; output heights are $O_n(L+h+p)$ and post-seed construction costs are polynomial in those quantities. This is not an exact membership claim for an unperturbed target outside the rational-factor set.

For a full matrix outside $C_q$, the cone is closed, so real separation applies. When the real fiber is feasible but rationally impossible, no real separating form can certify that arithmetic impossibility. For a nonclosed observation image, even real infeasibility can lack strict separation. The carrier theorem and the local quadratic-form obstruction give the correct alternatives rather than conflating these cases.

\subsection{Verification and the relation to prior local-pivot obstructions}
The accompanying standard-library scripts verify exact low-rank obstruction matrices, their range restrictions and finite modular obstructions; rational null and strict factors; an algebraic totally-real factor identity; signed Cayley identities; and homogeneous circuit blocks. A separate verifier imports no generator code and rejects deliberately corrupted data. These are checks of instances and interfaces. Neither a production general rational-point solver nor a general isotropy implementation is claimed.

The earlier Manuscript II branch proves exponential prefixes for every strictly descending extreme-ray adjacent rule on robust nonmembers of a fixed nonsimplicial column cone. The present theorem does not transfer that lower bound to every Lorentz cone, every initialization, or every algorithm. It instead answers a different generality question: the arithmetic content of a real boundary factorization depends on its quadratic zero domain. A short description and good interior conditioning do not alone control its rational feasibility, and rational certificate feasibility does not alone control a chosen local adjacency graph. The product reduction gives an independent reason why a proposed universal scalar condition for all matrix-cone problems would need qualifications.

\section{Priority boundaries and interpretation}
The real cone identity in Lemma~\ref{lem:balance} is not new: quadratic rank-one generation and semidefinite representations for one quadratic constraint have a substantial literature, including Sturm--Zhang \cite{SZ}. Nishijima--Louren\c{c}o \cite{NL} classify real faces of copositive and completely positive cones over a second-order cone. Their real facial classification is an important comparator, not something superseded merely by using arithmetic terminology.

Rational interior reconstruction has classical antecedents, including rational completely positive factorization \cite{DSV}; the present reserve proof is included to make the precise density hypothesis transparent. Rational orthogonal density, signed Cayley parametrizations, and quantitative rational approximation also predate this work \cite{Schmutz,LO,TWZ}. The exact width theorem is a consequence of these facts and positive-scaling openness, not a new parametrization of the orthogonal group. The $r+3$ unrestricted bound is derived from Meyer's theorem and is credited as classical quadratic-form arithmetic.

The Hasse--Minkowski principle, Meyer threshold, and small-zero bounds are classical \cite{CasselsBook,CasselsSmall,Dietmann}. The candidate new content is their exact integration into constrained rank-one generation and arbitrary affine Gram completion, including the maximal-rational-support step, the single-isotropy-test carrier, and the sharp rational spherical-moment obstruction. The product-cone circuit reduction supplies an explicit scope boundary. Burer--Dong \cite{BD} already represent nonconvex QCQPs through generalized completely positive cones over products of second-order cones and over PSD cones. Their bounded real convexification theorem and proof have been inspected; the present formulation isolates exact rational feasibility, only two homogeneous observations, unrestricted widths, and a binary-height lower bound without a boundedness promise. We do not claim that conic universality itself is new. We have not established that no earlier publication contains an equivalent formulation.

The earlier qualitative generalized-perfect-matrix theory \cite{OS} is also not a new result here. No blanket rational boundary factorization theorem is asserted for completely positive matrices or arbitrary rational cones. In particular, the Lorentz counterexamples here concern a different column cone and do not depend on reported 2026 orthant counterexamples.

\paragraph{Publication status.} The prior registered LAA comparison has not been completed or scored by this continuation. No numerical journal percentile, first-priority certificate, acceptance probability, external human verification, or proof-assistant verification is asserted. The contribution should be assessed primarily by the complete arithmetic classification and its affine and moment consequences, not by the number of examples or the name of a target venue.

\begin{thebibliography}{99}
\bibitem{SZ} J. F. Sturm and S. Zhang, On cones of nonnegative quadratic functions, \emph{Mathematics of Operations Research} 28 (2003), 246--267. DOI: 10.1287/moor.28.2.246.14485.
\bibitem{BD} S. Burer and H. Dong, Representing quadratically constrained quadratic programs as generalized copositive programs, \emph{Operations Research Letters} 40 (2012), 203--206. DOI: 10.1016/j.orl.2012.02.001.
\bibitem{NL} M. Nishijima and B. F. Louren\c{c}o, Facial structure of copositive and completely positive cones over a second-order cone, \emph{Linear Algebra and its Applications} 743 (2026), 211--249. DOI: 10.1016/j.laa.2026.04.019. Preprint arXiv:2502.04006v2.
\bibitem{SEDZ} M. Safey El Din and L. Zhi, Computing rational points in convex semi-algebraic sets and sum of squares decompositions, \emph{SIAM Journal on Optimization} 20 (2010), 2876--2889. DOI: 10.1137/090772459. In particular Theorem 1.1 and the real-algebraic sampling preliminaries.
\bibitem{CasselsBook} J. W. S. Cassels, \emph{Rational Quadratic Forms}, London Mathematical Society Monographs 13, Academic Press, 1978.
\bibitem{CasselsSmall} J. W. S. Cassels, Bounds for the least solutions of homogeneous quadratic equations, \emph{Proceedings of the Cambridge Philosophical Society} 51 (1955), 262--264.
\bibitem{Fukshansky} L. Fukshansky, Heights and quadratic forms: Cassels' theorem and its generalizations, \emph{Contemporary Mathematics} 587 (2013), 77--94. DOI: 10.1090/conm/587/11673.
\bibitem{Dietmann} R. Dietmann, Small zeros of quadratic forms avoiding a finite number of prescribed hyperplanes, \emph{Canadian Mathematical Bulletin} 52 (2009), 63--65. DOI: 10.4153/CMB-2009-007-7.
\bibitem{Schmutz} E. Schmutz, Rational points on the unit sphere, \emph{Central European Journal of Mathematics} 6 (2008), 482--487. DOI: 10.2478/s11533-008-0038-4.
\bibitem{LO} H. Liebeck and A. Osborn, The generation of all rational orthogonal matrices, \emph{American Mathematical Monthly} 98 (1991), 131--133.
\bibitem{TWZ} Q. Tang, W. Wang and H. Zhang, Generating all regular rational orthogonal matrices, \emph{Linear Algebra and its Applications} (2025). DOI: 10.1016/j.laa.2025.06.008. Preprint arXiv:2410.11246, whose Theorem 3 recalls Liebeck--Osborn.
\bibitem{DSV} M. Dutour Sikiri\'c, A. Sch\"urmann and F. Vallentin, Rational factorizations of completely positive matrices, \emph{Linear Algebra and its Applications} 523 (2017), 46--51. DOI: 10.1016/j.laa.2017.02.017.
\bibitem{OS} A. Oertel and A. Sch\"urmann, Generalized Perfect Matrices, arXiv:2602.05841v2 (2026).
\end{thebibliography}

\end{document}
